% Claim to Fame write-up (max. 1 page). Build: cd report && latexmk -pdf report.tex
% Numbers come from numbers.tex (python3 -I fetch.py), figures from plots.py. Lines marked CHECK need a team confirmation.
\documentclass[10pt]{article}
\usepackage[a4paper,margin=1.5cm,top=1.1cm,bottom=1.1cm]{geometry}
\usepackage[T1]{fontenc}
\usepackage{lmodern,microtype,graphicx,amsmath,enumitem,titlesec,xcolor,tabularx}
\usepackage[hidelinks]{hyperref}
% generated by fetch.py, do not edit
\newcommand{\FableFrom}{22}
\newcommand{\PreLast}{21}
\newcommand{\NGames}{100}
\newcommand{\NTeams}{12}
\newcommand{\Rank}{1}
\newcommand{\Total}{1,441k}
\newcommand{\BestOther}{1,174k}
\newcommand{\Lead}{267k}
\newcommand{\WorstRank}{9}
\newcommand{\WorstRankGame}{20}
\newcommand{\MinLead}{112k}
\newcommand{\MinLeadGame}{24}
\newcommand{\RankOneSince}{54}
\newcommand{\PreNet}{$-$12k}
\newcommand{\PreMean}{$-$0.6k}
\newcommand{\PostMean}{18.4k}
\newcommand{\PreWin}{48}
\newcommand{\PostWin}{87}
\newcommand{\PreTop}{10}
\newcommand{\PostTop}{62}
\newcommand{\PostLosses}{8}
\newcommand{\PostGames}{79}


\pagestyle{empty}
\setlength{\parindent}{0pt}
\setlength{\parskip}{2pt}
\titleformat{\section}{\normalsize\bfseries}{}{0pt}{}
\titlespacing*{\section}{0pt}{5pt}{1pt}
\setlist{nosep,leftmargin=1.2em,topsep=1pt}
\newcommand{\euro}{\,\texteuro}
\definecolor{c1}{HTML}{8E94A6}\definecolor{c2}{HTML}{B4BCF0}\definecolor{c3}{HTML}{7C86E6}
\definecolor{c4}{HTML}{3F49B5}\definecolor{c5}{HTML}{E0A93E}\definecolor{c6}{HTML}{D9478A}
\newcommand{\sq}[1]{\textcolor{#1}{\rule{1.6ex}{1.6ex}}}

\begin{document}

\begin{center}
  {\Large\bfseries Self-Learning Insurance Agent}\\[3pt]
  Kiebler Software Solutions \quad·\quad Tim Kiebler, Leo Driesch, Mattis Krauch, Sebastian Böhler % CHECK: names
\end{center}

\vspace{3pt}
{\small
\renewcommand{\arraystretch}{1.1}
\noindent\begin{tabularx}{\textwidth}{@{}l@{\;}X r r@{}}
  & \textbf{Our approach} & \textbf{Rounds} & \textbf{Cost per round}\\\hline
  \sq{c1} & Single agent (one model, low + medium effort) & 1--2 & \$0.004\\
  \sq{c2} & Ensemble, changing line-ups (8--20 model calls) & 3--6 & \$1.97\\
  \sq{c3} & Ensemble, final 8 members & 7--11 & \$0.86\\
  \sq{c4} & Ensemble with price references & 12--17 & \$0.86\\
  \sq{c6} & Multi-agent setup & 19, 21 & \$0.62\\
  \sq{c5} & \textbf{Self-learning persistent agent} (Fable 5.2, medium reasoning) & \textbf{18, 20, 22--100} & \textbf{\$0.77}\\\hline
\end{tabularx}\\[1pt]
\textbf{Table 1.} Which of our approaches submitted in each round, and its API cost per round (the agent: \$64.08 session cost at
API list prices over 83 rounds).}

\section{Method and uncertainty}
Our \textbf{Persistent Cost Agent} is \textbf{Claude Fable 5.2 (medium reasoning) in the Claude Code harness}:
one long-running session on a \texttt{/loop} plays every round and \textbf{keeps its context across games}
(1M-token window, \textbf{compacted} into a summary regularly at about 300--400k tokens), so it improves by
\textbf{in-context learning}. In its agent loop it reads and writes \textbf{Memory \& Tools}: files it writes itself on strategy and prices,
including \textbf{its own scripts to track prices and scores}, which it uses to verify quickly.
% CHECK: push these files to the repo and name 1-2 of them
For every invoice line, the agent uses its experience from earlier games and its world knowledge to decide
\textbf{first whether the line is covered and related}, \textbf{then estimates the hidden fair value $t$}, and from
that \textbf{chooses $a$ (our charge) and $b$ (the highest charge we accept) dynamically}. Uncertainty is weighed against the asymmetric payoff: wrongly rejecting a fair charge costs \textbf{1.5$\times$} it,
so the agent rejects only when confident, and charges slightly below its estimate of $t$. Python only checks the format and submits; \textbf{it does not change the numbers}.

\textbf{System 1, then System 2.} \textbf{System 1} is the agent's intuitive $(t,a,b)$ estimate, submitted after a median
of 13\,s. \textbf{System 2} refines it by reasoning (policy clauses, quantities, prices) and overwrites it, after a median
of 21\,s. A separate \textbf{Long-Horizon Evaluation Agent} reads the results of finished games and feeds
\textbf{adjustments} back to the cost agent.

\section{Evidence for our choices}
\begin{itemize}
  \item \textbf{Back-testing on all known cases}: approaches were run on every finished case and compared with the
        ranges for $t$ that the public results reveal. % CHECK: what exactly was back-tested, and by whom
  \item \textbf{Score from the challenge}: rank \Rank\ of \NTeams\ after \NGames\ rounds. Example, game 98: the policy listed the stolen watch on a valuables
        schedule; 8 of 11 teams rejected our fair 18,350\euro\ charge and paid 1.5$\times$. That line earned us 275k\euro.
  \item \textbf{Interaction with other agents}: each settlement shows how the other teams' agents charged and which
        charges they accepted; the agent uses this as feedback. % CHECK: meaning of this point
  \item \textbf{Several approaches run simultaneously}: the old ensemble kept deciding in parallel, without submitting.
\end{itemize}

\section{Resourcefulness}
Early on we used \textbf{token-heavy ensembles} (up to 20 model calls per round) to find good models. That was
expensive: it \textbf{burned through the OpenAI credits we were given} and would not have lasted 100 rounds. Its
benefit over a single model was also small, so we moved to a \textbf{single agent}, which did better in live play.
Multi-agent setups showed no advantage. The agent's whole tournament cost \textbf{169M tokens, 99\% of them cheap cache reads} of its persistent
context: about \textbf{\$64 at API list prices, \$0.77 per round} (the 8-model ensemble: \$0.86 per round).

\section{Reliability and automation}
The agent \textbf{runs fully autonomously} and \textbf{never missed a submission} in its 83 games. Median time from
release: \textbf{13\,s to the first estimate, 21\,s to the refinement} (slowest first estimate 38\,s; limit 60\,s). In games 19 and 21 a parallel multi-agent setup overwrote its answer; in game 20 that
setup crashed, and the agent's submission counted.

\section{Clarity and honesty}
The agent first has to build context, so it played \textbf{the first few rounds with reduced ability}.
\PostLosses\ of \PostGames\ games since the switch still lost money. We cannot fully separate the agent's effect from
the changing case mix, and an agent's decisions are harder to audit than a formula. We did not run ablations of System 1/2 or of the memory,
so their individual contributions are not measured.

\end{document}
